% ============================================================
% Lab Diagram - Device Functions (Reference), English
% Companion to assets/lab-diagram-tabloid-EN.pdf
% Compile: pdflatex 10-diagram-devices-EN.tex  (run twice)
% ============================================================
\documentclass[10pt,letterpaper]{article}

\usepackage[T1]{fontenc}
\usepackage[utf8]{inputenc}
\usepackage[letterpaper,margin=0.9in,top=0.8in,bottom=0.9in]{geometry}
\usepackage[scaled=.96]{helvet}
\renewcommand{\familydefault}{\sfdefault}
\usepackage{microtype}
\usepackage[table]{xcolor}
\usepackage{array}
\usepackage{longtable}
\usepackage{titlesec}
\usepackage{enumitem}
\usepackage{fancyhdr}
\PassOptionsToPackage{hyphens}{url}
\usepackage[colorlinks=true, urlcolor=PrimaryBlue, linkcolor=PrimaryBlue]{hyperref}

% ---- Palette (matches the diagram + pitch) -----------------
\definecolor{AccentBlue}{HTML}{3A83C2}
\definecolor{PrimaryBlue}{HTML}{1B5A8F}
\definecolor{Slate}{HTML}{273140}
\definecolor{Muted}{HTML}{5A6470}
\definecolor{DigiGreen}{HTML}{5B9E31}
\definecolor{Fase2Orange}{HTML}{D97706}
\definecolor{RFViolet}{HTML}{7A4FA0}

% ---- Layout ------------------------------------------------
\setlength{\parindent}{0pt}
\setlength{\parskip}{3.5pt}
\setcounter{secnumdepth}{0}
\titleformat{\section}
  {\large\bfseries\color{PrimaryBlue}}{}{0pt}{}
  [{\vspace{-2pt}\color{AccentBlue}\titlerule[1.1pt]}]
\titlespacing*{\section}{0pt}{10pt}{4pt}
\setlist[itemize]{leftmargin=1.2em, itemsep=2pt, topsep=2pt,
  label=\textcolor{AccentBlue}{\textbullet}}

\pagestyle{fancy}
\fancyhf{}
\renewcommand{\headrulewidth}{0pt}
\fancyfoot[L]{\footnotesize\color{Muted} Lab Diagram --- Device Functions \textbullet\ source: assets/lab-diagram-tabloid-EN.tex}
\fancyfoot[R]{\footnotesize\color{Muted} Page \thepage}

% ---- Helpers -----------------------------------------------
\newcommand{\kf}[1]{\textbf{\textcolor{PrimaryBlue}{#1}}}
\newcommand{\swatch}[1]{\textcolor{#1}{\rule[1pt]{16pt}{3pt}}}
\newcommand{\notebox}[1]{%
  \setlength{\fboxsep}{9pt}\setlength{\fboxrule}{0.8pt}%
  \noindent\fcolorbox{Fase2Orange}{Fase2Orange!7}{%
  \begin{minipage}{\dimexpr\linewidth-2\fboxsep-2\fboxrule\relax}#1\end{minipage}}}

% device tables: 2 columns, header repeats on page break
\newcommand{\dtophead}{%
  \rowcolor{PrimaryBlue}{\color{white}\bfseries Device} & {\color{white}\bfseries Function}\\}
\newenvironment{devtable}{%
  \setlength{\arrayrulewidth}{0.4pt}\arrayrulecolor{black!15}%
  \renewcommand{\arraystretch}{1.28}%
  \begin{longtable}{@{}>{\raggedright\arraybackslash}p{4.4cm}@{\hspace{9pt}}>{\raggedright\arraybackslash}p{11.9cm}@{}}
  \dtophead \hline \endfirsthead \dtophead \hline \endhead}
  {\end{longtable}}

\begin{document}

% ---- Title block -------------------------------------------
{\color{AccentBlue}\rule{\linewidth}{4pt}}
\vspace{7pt}
{\LARGE\bfseries\color{Slate} Lab Diagram --- Device Functions\par}
\vspace{3pt}
{\large\bfseries\color{PrimaryBlue} Reference companion to the tabloid diagram\par}
\vspace{3pt}
{\small\color{Muted} Starlink-Primary Resilient Gateway Reference Lab \textbullet\ Digi International \textbullet\ July 2026\par}

\vspace{6pt}
Companion to \href{run:assets/lab-diagram-tabloid-EN.pdf}{\texttt{assets/lab-diagram-tabloid-EN.pdf}}. It lists \kf{what each device in the diagram does}. This document does \textbf{not} modify the diagram; it only explains it.

\vspace{2pt}
{\small\color{Muted}\textit{Convention: one site in the sprint (Ring 0 = Site A); three sites + central control in the expansion (Rings 1--2). All base hardware is identical between sites; only the profile (native / bonding / candidate) rotates between units.}}

% ---- Central control ---------------------------------------
\section{Central Control --- cloud + lab host (Proxmox) \textbullet\ AS65000}
\begin{devtable}
\textbf{Digi Remote Manager Premier} (cloud) & Single management plane: golden config + per-site settings, alerts (Device Offline / noncompliance), Push Monitor webhook $\rightarrow$ SIEM, Digi Containers deployment, and CCCS OTA for the embedded sensor. \\ \hline
\textbf{Core / route collector} & Public overlay endpoint (IKEv2/IPsec) and eBGP peer for every site (AS65000). Collects the network each site originates. \\ \hline
\textbf{SIEM} & Receives Push Monitor webhooks and the MP255 agent stream; event / security correlation. \\ \hline
\textbf{Impairment platform} & Injects controlled loss / jitter / latency to reproduce Starlink conditions during tests. \\ \hline
\textbf{WAN Bonding server} & External Tier 3 (1 Gb/s) bonding server required by the WAN-bonding experiment (A/B vs SureLink). \\ \hline
\textbf{Lighthouse Enhance VM} \textit{(Phase 2)} & Opengear OOB management plane. \kf{Only needed if the OOB device is the Opengear OM1304}; with a \kf{Digi Connect IT 4} the OOB is managed by Digi Remote Manager instead (no separate VM). \\
\end{devtable}

% ---- Site A ------------------------------------------------
\section{Site A --- AS65001 \textbullet\ NATIVE profile \textbullet\ Ring 0 (2-week sprint)}
\begin{devtable}
\textbf{Starlink Performance Kit (Gen 3)} & Primary WAN. Copper to IX40 \texttt{WAN/ETH1}; delivers CGNAT \texttt{100.64.0.0/10} (public IPv4 optional). \\ \hline
\textbf{Carrier A --- 5G} & In-band backup WAN through the IX40 5G modem; own SIM/APN. \\ \hline
\textbf{Digi IX40-05} & The resilient gateway. Runs SureLink failover, the IKEv2/IPsec overlay, eBGP (AS65001) and the Digi Container. \\ \hline
\textbf{Digi Container} & Metrics collector running \textbf{on} the IX40 under resource limits (proves Containers without degrading forwarding). \\ \hline
\textbf{Service fabric (switch)} & Service network behind \texttt{SFP/ETH5} (1 Gb/s SFP, never SFP+): IPv4 \texttt{/24} + IPv6 \texttt{/64} advertised via eBGP. \\ \hline
\textbf{Digi XBee Hive} & RF gateway; isolates the sensor application from packet forwarding. \\ \hline
\textbf{Digi XBee 3 (Zigbee)} & Zigbee radio linking the Hive to the physical sensor. \\ \hline
\textbf{Physical sensor} & Measures temperature / door state / a digital input; not part of critical forwarding. \\ \hline
\textbf{Digi ConnectCore MP255} & Embedded SoM. DEY (Yocto) image; agent $\rightarrow$ SIEM (offline buffer + replay); OTA via CCCS. \\ \hline
\textbf{OOB console server} \textit{(Phase 2)} --- \textbf{Opengear OM1304-4E-C-LSP} and/or \textbf{Digi Connect IT 4} & Out-of-band rescue: serial console (DB9, \texttt{Login} mode) to the IX40, independent \kf{5G Carrier B}, management switch and independent power. Two interchangeable options --- see the note below. Candidate to deploy at \kf{Site B}. \\ \hline
\textbf{Carrier B --- 5G} & Independent carrier feeding the OOB path (different operator and power from Carrier A / Starlink). \\ \hline
\textbf{Primary PDU/UPS} & Powers Starlink + IX40 + load. \\ \hline
\textbf{Secondary PDU/UPS} & Powers the OOB device + modem B (independent power domain). \\
\end{devtable}

% ---- Sites B & C -------------------------------------------
\section{Sites B \& C (expansion \textbullet\ Rings 1--2)}
\begin{devtable}
\textbf{Site B --- AS65002} & Identical hardware to Site A. \kf{BONDING} profile: Starlink + 5G inside the bonded tunnel. Planned \kf{OOB via Digi Connect IT 4}. \\ \hline
\textbf{Site C --- AS65003} & Identical to Site A. \kf{CANDIDATE / CHAOS} profile: candidate containers, config and firmware; A/B/C profile rotation. \\
\end{devtable}

% ---- OOB options -------------------------------------------
\section{OOB console server --- two interchangeable options}
Both do the same job in the diagram (serial console to the IX40 + independent 5G + own power). They differ mainly in the management plane:

\vspace{4pt}
{\setlength{\arrayrulewidth}{0.4pt}\arrayrulecolor{black!15}\renewcommand{\arraystretch}{1.28}%
\begin{tabular}{@{}>{\raggedright\arraybackslash}p{3.3cm}@{\hspace{8pt}}>{\raggedright\arraybackslash}p{6.4cm}@{\hspace{8pt}}>{\raggedright\arraybackslash}p{6.4cm}@{}}
\rowcolor{PrimaryBlue}{\color{white}\bfseries\ } & {\color{white}\bfseries Opengear OM1304-4E-C-LSP} & {\color{white}\bfseries Digi Connect IT 4}\\ \hline
\textbf{Serial to IX40} & DB9 $\leftrightarrow$ RJ45, adapter \textbf{319015} (Opengear X2 pinout) & RJ45 RS-232, Digi/Cisco-pinout cable \\ \hline
\textbf{Cellular} & 5G Carrier B & LTE Cat 4 Carrier B (Digi CORE modem) \\ \hline
\textbf{Mgmt ports} & 4-port mgmt switch & 2$\times$ Ethernet 10/100 (add a small switch for 4 ports) \\ \hline
\textbf{Managed by} & Opengear \textbf{Lighthouse} (needs the Phase-2 VM) & \textbf{Digi Remote Manager} (same tenant as the routers $\rightarrow$ one plane) \\
\end{tabular}}

\vspace{8pt}
\notebox{\kf{Under consideration: Connect IT 4 at Site B.} Choosing it unifies management under DRM and removes the Lighthouse VM; keeping the OM1304 preserves vendor-diverse OOB. Both are listed on purpose --- the diagram is left unchanged for now.}

% ---- Link legend -------------------------------------------
\section{Link legend (as drawn)}
\begin{itemize}
  \item \swatch{PrimaryBlue}\ \ Starlink primary WAN / overlay + eBGP (thick blue)
  \item \swatch{AccentBlue}\ \ 5G in-band backup / inter-site links
  \item \swatch{DigiGreen}\ \ Service fabric (1 Gb/s SFP fiber)
  \item \swatch{Muted}\ \ Management / LAN --- HTTPS, API, DRM (dotted)
  \item \swatch{Fase2Orange}\ \ Out-of-band / Phase 2 elements (dashed)
  \item \swatch{RFViolet}\ \ Zigbee / cellular RF
  \item Dashed box = Phase 2
\end{itemize}

\vspace{4pt}
{\footnotesize\color{Muted} Source: \texttt{assets/lab-diagram-tabloid-EN.tex} and \texttt{03-arquitectura.md}. Product logos are property of their respective owners.}

\end{document}
